\section*{NeurIPS Paper Checklist}

\begin{enumerate}

\item {\bf Claims}
    \item[] Question: Do the main claims made in the abstract and introduction accurately reflect the paper's contributions and scope?
    \item[] Answer: \answerYes{}
    \item[] Justification: The abstract and introduction state the main claims as a cross-architecture empirical audit: significant Method$\times$Domain interaction, unstable cross-domain method transfer, and architecture-specific efficiency tradeoffs. Representation-similarity analysis is explicitly framed as preliminary diagnostic evidence rather than as a primary claim.

\item {\bf Limitations}
    \item[] Question: Does the paper discuss the limitations of the work performed by the authors?
    \item[] Answer: \answerYes{}
    \item[] Justification: Limitations are discussed in Section~5 under a dedicated paragraph. These include limited domain coverage, uneven seed coverage across methods, the preliminary status of CKA diagnostics, the treatment of divergence outliers, partial MASE coverage, and the small $n{=}12$ LODO evaluation grid.

\item {\bf Theory assumptions and proofs}
    \item[] Question: For each theoretical result, does the paper provide the full set of assumptions and a complete (and correct) proof?
    \item[] Answer: \answerNA{}
    \item[] Justification: The paper does not present formal theoretical results, theorems, or proofs; it is an empirical benchmarking and analysis study.

\item {\bf Experimental result reproducibility}
    \item[] Question: Does the paper fully disclose all the information needed to reproduce the main experimental results of the paper to the extent that it affects the main claims and/or conclusions of the paper (regardless of whether the code and data are provided or not)?
    \item[] Answer: \answerYes{}
    \item[] Justification: Section~3 and the appendix reproducibility statement specify datasets, model families, adaptation methods, ranks, loci, statistical tests, compute setup, planned seeds, completed-run accounting, software versions, and Hydra-based training entry points needed to reproduce the reported analyses.

\item {\bf Open access to data and code}
    \item[] Question: Does the paper provide open access to the data and code, with sufficient instructions to faithfully reproduce the main experimental results, as described in supplemental material?
    \item[] Answer: \answerYes{}
    \item[] Justification: Code, configs, run manifest, headline-number artifacts, and the Croissant 1.0 manifest are released for review via an anonymous URL (\texttt{https://anonymous.4open.science/r/tsfm-peft-bench-anon}), and the supplementary tarball bundles the same artifacts plus per-run JSONs. The Reproducibility Statement (Appendix~\ref{sec:repro}) and Datasheet (Appendix~\ref{sec:datasheet}) describe the exact reproduction commands. Source datasets are publicly available; the README documents per-platform hosting plans for camera-ready release.

\item {\bf Experimental setting/details}
    \item[] Question: Does the paper specify all the training and test details (e.g., data splits, hyperparameters, how they were chosen, type of optimizer) necessary to understand the results?
    \item[] Answer: \answerYes{}
    \item[] Justification: Section~3 describes datasets, adaptation axes, experiment scale, and the primary statistical methodology. The appendix reproducibility statement adds hardware, software versions, the planned seed set, completed observation counts, and the Hydra entry points used to launch training.

\item {\bf Experiment statistical significance}
    \item[] Question: Does the paper report error bars suitably and correctly defined or other appropriate information about the statistical significance of the experiments?
    \item[] Answer: \answerYes{}
    \item[] Justification: The main claims are supported by ANOVA $F$-statistics and $\eta^2$ effect sizes, mean$\pm$std MAE in the efficiency table, and Wilson confidence intervals for the LODO transfer experiment. The appendix additionally reports ANOVA degrees of freedom and exploratory robustness analyses.

\item {\bf Experiments compute resources}
    \item[] Question: For each experiment, does the paper provide sufficient information on the computer resources (type of compute workers, memory, time of execution) needed to reproduce the experiments?
    \item[] Answer: \answerYes{}
    \item[] Justification: The appendix reproducibility statement reports the hardware (4$\times$ RTX 3090, 24\,GB each), precision (FP16/BF16), approximate total compute (${\sim}$2,250 GPU-hours across 885 completed runs), and software stack. Table~\ref{tab:efficiency} reports per-run wall-clock time by method.

\item {\bf Code of ethics}
    \item[] Question: Does the research conducted in the paper conform, in every respect, with the NeurIPS Code of Ethics \url{https://neurips.cc/public/EthicsGuidelines}?
    \item[] Answer: \answerYes{}
    \item[] Justification: The work is an empirical analysis of public datasets and pre-trained models, with no human subjects or sensitive data collection. The submission preserves anonymity and discusses failure modes and deployment caution in the appendix broader impacts section.

\item {\bf Broader impacts}
    \item[] Question: Does the paper discuss both potential positive societal impacts and negative societal impacts of the work performed?
    \item[] Answer: \answerYes{}
    \item[] Justification: The appendix broader impacts section discusses positive impacts (more reliable adaptation, reduced wasted compute) and negative impacts (overconfident transfer under domain shift, potential surveillance misuse), together with the mitigation recommendation of domain-specific validation before deployment.

\item {\bf Safeguards}
    \item[] Question: Does the paper describe safeguards that have been put in place for responsible release of data or models that have a high risk for misuse (e.g., pre-trained language models, image generators, or scraped datasets)?
    \item[] Answer: \answerNA{}
    \item[] Justification: The paper does not release a new model, scraped dataset, or dual-use asset. It analyzes adaptation behavior of existing public models and datasets.

\item {\bf Licenses for existing assets}
    \item[] Question: Are the creators or original owners of assets (e.g., code, data, models), used in the paper, properly credited and are the license and terms of use explicitly mentioned and properly respected?
    \item[] Answer: \answerYes{}
    \item[] Justification: The appendix enumerates verified upstream licenses (Apache-2.0 for Chronos/TimesFM, MIT for MOMENT, CC-BY-NC-4.0 for Moirai, CC-BY-ND-4.0 for ETTm1). The Exchange Rate benchmark and the bundled SMD code are clearly credited to their upstream maintainers (\texttt{laiguokun/multivariate-time-series-data} and \texttt{NetManAIOps/OmniAnomaly}); we use them under the upstream repositories' code licenses (MIT) and do not redistribute raw source data. Our own contributions (code, manifest, Croissant file) are released under Apache-2.0. Asset-license caveats for the two upstream sources without standalone dataset license files are documented transparently in the Reproducibility Statement.

\item {\bf New assets}
    \item[] Question: Are new assets introduced in the paper well documented and is the documentation provided alongside the assets?
    \item[] Answer: \answerYes{}
    \item[] Justification: We introduce \textbf{TSFM-PEFT-Bench}, a benchmark of 882 paper-included primary-model PEFT runs (972 total). Documentation includes: a Croissant 1.0 manifest with mandatory Responsible-AI metadata (\texttt{tsfm\_peft\_bench.croissant.json}); a datasheet following~\citet{gebru2021datasheets} (Appendix~\ref{sec:datasheet}); per-domain statistics (Appendix~\ref{sec:domain-stats}); hyperparameters (Appendix~\ref{sec:hyperparams}); per-cell raw MAE (Appendix~\ref{sec:per-cell-mae}); trainable parameter counts (Appendix~\ref{sec:params}); reproducibility statement (Appendix~\ref{sec:repro}); and an \texttt{LICENSE} file (Apache-2.0). Hosting plan is documented in the README.

\item {\bf Crowdsourcing and research with human subjects}
    \item[] Question: For crowdsourcing experiments and research with human subjects, does the paper include the full text of instructions given to participants and screenshots, if applicable, as well as details about compensation (if any)?
    \item[] Answer: \answerNA{}
    \item[] Justification: The research does not involve crowdsourcing or experiments with human subjects.

\item {\bf Institutional review board (IRB) approvals or equivalent for research with human subjects}
    \item[] Question: Does the paper describe potential risks incurred by study participants, whether such risks were disclosed to the subjects, and whether Institutional Review Board (IRB) approvals (or an equivalent approval/review based on the requirements of your country or institution) were obtained?
    \item[] Answer: \answerNA{}
    \item[] Justification: The research does not involve human subjects, so IRB approval is not applicable.

\item {\bf Declaration of LLM usage}
    \item[] Question: Does the paper describe the usage of LLMs if it is an important, original, or non-standard component of the core methods in this research? Note that if the LLM is used only for writing, editing, or formatting purposes and does \emph{not} impact the core methodology, scientific rigor, or originality of the research, declaration is not required.
    \item[] Answer: \answerNA{}
    \item[] Justification: LLMs were not used as part of the core methodology, experiments, or scientific claims. For transparency, the Reproducibility Statement (Appendix~\ref{sec:repro}) records that Anthropic's Claude and OpenAI's Codex assisted only with English drafting/LaTeX editing, academic-style polishing, and Python code review/bug-patch assistance; these tools did not design experiments, generate empirical results, choose baselines, or alter any reported number. Per the NeurIPS 2026 LLM policy, this usage is exempt from disclosure; we declare it nonetheless.

\end{enumerate}
